\documentclass[10pt,a4paper]{amsart}
\usepackage{amsmath,amssymb,amsthm,graphicx}
% Retained minimal preamble for the Stacks Project chapter more-algebra.tex (More on Algebra).
% The chapter file pinned-original-source.tex carries only \input{preamble} on its first line, so no
% preamble text lives inside the chapter itself.  The Stacks Project uses one single shared preamble for
% every chapter; the snapshot stages it identically as sources/stacks-chow/preamble.tex and
% sources/stacks-etale-cohomology/preamble.tex.  The block below is copied mechanically, line for line,
% from sources/stacks-chow/preamble.tex lines 14-18 (the Xy-pic setup) and 154-217 (the native theorem
% declarations together with the mathematical macros).  Nothing was added, removed or redefined; the
% wrapper class and the amsmath/amssymb/amsthm/graphicx base packages are supplied by the assembler.
% Chapter: more-algebra.tex -- More on Algebra.
% For commutative diagrams we use Xy-pic
\usepackage[all]{xy}

% We use 2cell for 2-commutative diagrams.
\xyoption{2cell}

% Theorem environments.
%
\theoremstyle{plain}
\newtheorem{theorem}[subsection]{Theorem}
\newtheorem{proposition}[subsection]{Proposition}
\newtheorem{lemma}[subsection]{Lemma}

\theoremstyle{definition}
\newtheorem{definition}[subsection]{Definition}
\newtheorem{example}[subsection]{Example}
\newtheorem{exercise}[subsection]{Exercise}
\newtheorem{situation}[subsection]{Situation}

\theoremstyle{remark}
\newtheorem{remark}[subsection]{Remark}
\newtheorem{remarks}[subsection]{Remarks}

\numberwithin{equation}{subsection}

% Macros
%
\def\lim{\mathop{\mathrm{lim}}\nolimits}
\def\colim{\mathop{\mathrm{colim}}\nolimits}
\def\Spec{\mathop{\mathrm{Spec}}}
\def\Hom{\mathop{\mathrm{Hom}}\nolimits}
\def\Ext{\mathop{\mathrm{Ext}}\nolimits}
\def\SheafHom{\mathop{\mathcal{H}\!\mathit{om}}\nolimits}
\def\SheafExt{\mathop{\mathcal{E}\!\mathit{xt}}\nolimits}
\def\Sch{\mathit{Sch}}
\def\Mor{\mathop{\mathrm{Mor}}\nolimits}
\def\Ob{\mathop{\mathrm{Ob}}\nolimits}
\def\Sh{\mathop{\mathit{Sh}}\nolimits}
\def\NL{\mathop{N\!L}\nolimits}
\def\CH{\mathop{\mathrm{CH}}\nolimits}
\def\proetale{{pro\text{-}\acute{e}tale}}
\def\etale{{\acute{e}tale}}
\def\QCoh{\mathit{QCoh}}
\def\Ker{\mathop{\mathrm{Ker}}}
\def\Im{\mathop{\mathrm{Im}}}
\def\Coker{\mathop{\mathrm{Coker}}}
\def\Coim{\mathop{\mathrm{Coim}}}

% Boxtimes
%
\DeclareMathSymbol{\boxtimes}{\mathbin}{AMSa}{"02}

%
% Macros for moduli stacks/spaces
%
\def\QCohstack{\mathcal{QC}\!\mathit{oh}}
\def\Cohstack{\mathcal{C}\!\mathit{oh}}
\def\Spacesstack{\mathcal{S}\!\mathit{paces}}
\def\Quotfunctor{\mathrm{Quot}}
\def\Hilbfunctor{\mathrm{Hilb}}
\def\Curvesstack{\mathcal{C}\!\mathit{urves}}
\def\Polarizedstack{\mathcal{P}\!\mathit{olarized}}
\def\Complexesstack{\mathcal{C}\!\mathit{omplexes}}
% \Pic is the operator that assigns to X its picard group, usage \Pic(X)
% \Picardstack_{X/B} denotes the Picard stack of X over B
% \Picardfunctor_{X/B} denotes the Picard functor of X over B
\def\Pic{\mathop{\mathrm{Pic}}\nolimits}
\def\Picardstack{\mathcal{P}\!\mathit{ic}}
\def\Picardfunctor{\mathrm{Pic}}
\def\Deformationcategory{\mathcal{D}\!\mathit{ef}}

% --- appended 2026-10-12 (Chief): the project own definitions of the reference/slogan/history
% display environments, copied mechanically from the shared Stacks preamble lines 8-12
% (\usepackage{verbatim} plus the three \newenvironment lines, which the chapter preamble keeps
% verbatim from the project preamble). Nothing added, removed or redefined: these environments
% carry no mathematical content and the project itself typesets them as comments. Without these
% three lines a counted statement whose own text uses them fails with Environment slogan undefined.
\usepackage{verbatim}
\newenvironment{reference}{\comment}{\endcomment}
\newenvironment{slogan}{\comment}{\endcomment}
\newenvironment{history}{\comment}{\endcomment}
\usepackage{cleveref}
\crefname{originallemma}{original lemma}{original lemmas}
\Crefname{originallemma}{Original lemma}{Original lemmas}
\makeatletter
\newlabel{algebra-lemma-characterize-dvr}{{00PD}{source}{}{sourcepointer-1}{}}
\newlabel{algebra-lemma-characterize-dvr@cref}{{[originallemma][0][]00PD}{[1][0][]source}{}{}{}}
\makeatother
\crefname{originallemma}{original lemma}{original lemmas}
\Crefname{originallemma}{Original lemma}{Original lemmas}
\makeatletter
\newlabel{algebra-lemma-criterion-normal}{{031S}{source}{}{sourcepointer-2}{}}
\newlabel{algebra-lemma-criterion-normal@cref}{{[originallemma][0][]031S}{[1][0][]source}{}{}{}}
\makeatother
\crefname{originallemma}{original lemma}{original lemmas}
\Crefname{originallemma}{Original lemma}{Original lemmas}
\makeatletter
\newlabel{algebra-lemma-length-independent}{{00IX}{source}{}{sourcepointer-3}{}}
\newlabel{algebra-lemma-length-independent@cref}{{[originallemma][0][]00IX}{[1][0][]source}{}{}{}}
\makeatother
\crefname{originallemma}{original lemma}{original lemmas}
\Crefname{originallemma}{Original lemma}{Original lemmas}
\makeatletter
\newlabel{algebra-lemma-minimal-over-1}{{00KV}{source}{}{sourcepointer-4}{}}
\newlabel{algebra-lemma-minimal-over-1@cref}{{[originallemma][0][]00KV}{[1][0][]source}{}{}{}}
\makeatother
\crefname{originallemma}{original lemma}{original lemmas}
\Crefname{originallemma}{Original lemma}{Original lemmas}
\makeatletter
\newlabel{algebra-lemma-normal-domain-intersection-localizations-height-1}{{031T}{source}{}{sourcepointer-5}{}}
\newlabel{algebra-lemma-normal-domain-intersection-localizations-height-1@cref}{{[originallemma][0][]031T}{[1][0][]source}{}{}{}}
\makeatother
\crefname{originallemma}{original lemma}{original lemmas}
\Crefname{originallemma}{Original lemma}{Original lemmas}
\makeatletter
\newlabel{algebra-lemma-silly}{{00DS}{source}{}{sourcepointer-6}{}}
\newlabel{algebra-lemma-silly@cref}{{[originallemma][0][]00DS}{[1][0][]source}{}{}{}}
\makeatother
\crefname{originallemma}{original lemma}{original lemmas}
\Crefname{originallemma}{Original lemma}{Original lemmas}
\makeatletter
\newlabel{algebra-lemma-zero-at-ass-zero}{{0311}{source}{}{sourcepointer-7}{}}
\newlabel{algebra-lemma-zero-at-ass-zero@cref}{{[originallemma][0][]0311}{[1][0][]source}{}{}{}}
\makeatother
\crefname{originallemma}{original lemma}{original lemmas}
\Crefname{originallemma}{Original lemma}{Original lemmas}
\makeatletter
\newlabel{lemma-minprimespoly}{{38016}{source}{}{sourcepointer-8}{}}
\newlabel{lemma-minprimespoly@cref}{{[originallemma][0][]38016}{[1][0][]source}{}{}{}}
\makeatother
\crefname{originallemma}{original lemma}{original lemmas}
\Crefname{originallemma}{Original lemma}{Original lemmas}
\makeatletter
\newlabel{lemma-minprimespolyhigher}{{38297}{source}{}{sourcepointer-9}{}}
\newlabel{lemma-minprimespolyhigher@cref}{{[originallemma][0][]38297}{[1][0][]source}{}{}{}}
\makeatother
\crefname{originallemma}{original lemma}{original lemmas}
\Crefname{originallemma}{Original lemma}{Original lemmas}
\makeatletter
\newlabel{lemma-radical-element}{{38138}{source}{}{sourcepointer-10}{}}
\newlabel{lemma-radical-element@cref}{{[originallemma][0][]38138}{[1][0][]source}{}{}{}}
\makeatother
\crefname{originallemma}{original lemma}{original lemmas}
\Crefname{originallemma}{Original lemma}{Original lemmas}
\makeatletter
\newlabel{lemma-syspar}{{38108}{source}{}{sourcepointer-11}{}}
\newlabel{lemma-syspar@cref}{{[originallemma][0][]38108}{[1][0][]source}{}{}{}}
\makeatother
\crefname{originallemma}{original lemma}{original lemmas}
\Crefname{originallemma}{Original lemma}{Original lemmas}
\makeatletter
\newlabel{lemma-sysparhigher}{{38348}{source}{}{sourcepointer-12}{}}
\newlabel{lemma-sysparhigher@cref}{{[originallemma][0][]38348}{[1][0][]source}{}{}{}}
\makeatother
\title{Over a catenary Noetherian local normal domain every radical ideal contains a nonzero element whose quotient is reduced}
\author{The Stacks Project Authors}
\date{stacks-project:more-algebra.tex:local-snapshot-20261012; GNU FDL 1.2 or later}
\begin{document}
\maketitle
\setcounter{subsection}{1}
\noindent\textit{Editorial note.} All delivered mathematical text is the author's own, copied from the pinned Stacks Project chapter file more-algebra.tex; no statement or proof step was rewritten, reordered or removed.\par\medskip
\noindent\textit{Editorial note.} The chapter file pinned-original-source.tex carries only \textbackslash{}input\{preamble\} on its first line, so no preamble text lives inside the chapter. The retained file sources/stacks-more-algebra/retained-preamble.tex is a mechanical, line-for-line copy of the project's single shared preamble as staged in the snapshot, together with the project's own definitions of the reference/slogan/history display environments, which the project itself typesets as comments; nothing was added, removed or redefined by this card and no package was added.\par\medskip
\noindent\textit{Editorial note.} The source URL is the source-level Stacks Project URL, as required when more than one card is produced from one source; the per-claim tag is carried by claim\_key.\par\medskip
\noindent\textit{Editorial note.} The counted statement span includes the author's own \textbackslash{}begin\{reference\} block, in which the author notes that the cited work 'has this result without the assumption that the ring is catenary' and cites it; that block is delivered verbatim. Because the project defines the reference environment as a comment environment, the citation is typeset as a comment and contributes no unresolved citation. The author's typos in the counted proof span ('q\_1 cap ... cap q\_s' for one occurrence of the localizations, and the index j' used in the display defining c) are reproduced unchanged.\par\medskip
\noindent\textit{Allowed standard prerequisites (editorial).}\begin{itemize}
\item Catenary local rings: every height one prime has dimension d - 1, and normality giving R = intersection of the localizations at height one primes. used at proof lines 38405-38407 and invoked through the cited Algebra lemmas at proof lines 38432-38434. Makes Spec(R/(f + h)) equidimensional and lets reducedness be tested on the height one primes.
\item Discrete valuation rings at height one primes of a normal domain and the divisor of an element. used at proof lines 38414-38426. Supplies the valuations ord\_\{q\_j\} whose values m\_j are exactly what must be shown to be 1.
\item The height-one-polynomial bound on the number of minimal primes. invoked at proof lines 38437-38438 and applied at proof lines 38446-38455. Is the input that bounds t(f + h) uniformly, enabling the maximality argument.
\item Prime avoidance and the second symbolic power of a height one prime. used at proof lines 38487-38500. Produces the elements a\_j and b\_j from which the correcting element c is assembled.
\end{itemize}\medskip
\ifdefined\hypertarget\hypertarget{sourcepointer-1}{}\fi
\noindent\textit{Source reference.} Original lemma 00PD: focused-seeds/AGP/sources/stacks-snapshot-20261011/tags.txt, source lines 927--927. Cross-chapter navigation pointer into the Stacks Project; the standard fact it supplies is applied inside the counted proof span and the argument resting on it stays there.\par\medskip
\ifdefined\hypertarget\hypertarget{sourcepointer-2}{}\fi
\noindent\textit{Source reference.} Original lemma 031S: focused-seeds/AGP/sources/stacks-snapshot-20261011/tags.txt, source lines 3995--3995. Cross-chapter navigation pointer into the Stacks Project; the standard fact it supplies is applied inside the counted proof span and the argument resting on it stays there.\par\medskip
\ifdefined\hypertarget\hypertarget{sourcepointer-3}{}\fi
\noindent\textit{Source reference.} Original lemma 00IX: focused-seeds/AGP/sources/stacks-snapshot-20261011/tags.txt, source lines 695--695. Cross-chapter navigation pointer into the Stacks Project; the standard fact it supplies is applied inside the counted proof span and the argument resting on it stays there.\par\medskip
\ifdefined\hypertarget\hypertarget{sourcepointer-4}{}\fi
\noindent\textit{Source reference.} Original lemma 00KV: focused-seeds/AGP/sources/stacks-snapshot-20261011/tags.txt, source lines 765--765. Cross-chapter navigation pointer into the Stacks Project; the standard fact it supplies is applied inside the counted proof span and the argument resting on it stays there.\par\medskip
\ifdefined\hypertarget\hypertarget{sourcepointer-5}{}\fi
\noindent\textit{Source reference.} Original lemma 031T: focused-seeds/AGP/sources/stacks-snapshot-20261011/tags.txt, source lines 3996--3996. Cross-chapter navigation pointer into the Stacks Project; the standard fact it supplies is applied inside the counted proof span and the argument resting on it stays there.\par\medskip
\ifdefined\hypertarget\hypertarget{sourcepointer-6}{}\fi
\noindent\textit{Source reference.} Original lemma 00DS: focused-seeds/AGP/sources/stacks-snapshot-20261011/tags.txt, source lines 508--508. Cross-chapter navigation pointer into the Stacks Project; the standard fact it supplies is applied inside the counted proof span and the argument resting on it stays there.\par\medskip
\ifdefined\hypertarget\hypertarget{sourcepointer-7}{}\fi
\noindent\textit{Source reference.} Original lemma 0311: focused-seeds/AGP/sources/stacks-snapshot-20261011/tags.txt, source lines 3968--3968. Cross-chapter navigation pointer into the Stacks Project; the standard fact it supplies is applied inside the counted proof span and the argument resting on it stays there.\par\medskip
\ifdefined\hypertarget\hypertarget{sourcepointer-8}{}\fi
\noindent\textit{Source reference.} Original lemma 38016: sources/stacks-more-algebra/pinned-original-source.tex, source lines 38016--38016. Same-chapter navigation pointer; the statement it names is standard for the section and the argument that consumes it stays inside the counted proof span.\par\medskip
\ifdefined\hypertarget\hypertarget{sourcepointer-9}{}\fi
\noindent\textit{Source reference.} Original lemma 38297: sources/stacks-more-algebra/pinned-original-source.tex, source lines 38297--38297. Same-chapter navigation pointer; the statement it names is standard for the section and the argument that consumes it stays inside the counted proof span.\par\medskip
\ifdefined\hypertarget\hypertarget{sourcepointer-10}{}\fi
\noindent\textit{Source reference.} Original lemma 38138: sources/stacks-more-algebra/pinned-original-source.tex, source lines 38138--38138. Same-chapter navigation pointer; the statement it names is standard for the section and the argument that consumes it stays inside the counted proof span.\par\medskip
\ifdefined\hypertarget\hypertarget{sourcepointer-11}{}\fi
\noindent\textit{Source reference.} Original lemma 38108: sources/stacks-more-algebra/pinned-original-source.tex, source lines 38108--38108. Same-chapter navigation pointer; the statement it names is standard for the section and the argument that consumes it stays inside the counted proof span.\par\medskip
\ifdefined\hypertarget\hypertarget{sourcepointer-12}{}\fi
\noindent\textit{Source reference.} Original lemma 38348: sources/stacks-more-algebra/pinned-original-source.tex, source lines 38348--38348. Same-chapter navigation pointer; the statement it names is standard for the section and the argument that consumes it stays inside the counted proof span.\par\medskip
\begin{proposition}
\label{proposition-propdimd}
\begin{reference}
\cite[Lemma 3.14]{Artin-Lipman} has this result without the
assumption that the ring is catenary
\end{reference}
Let $R$ be a catenary Noetherian local normal domain.
Let $J \subset R$ be a radical ideal.
Then there exists a nonzero element $f \in J$
such that $R/fR$ is reduced.
\end{proposition}

\begin{proof}
The proof is the same as that of
Lemma \ref{lemma-radical-element},
using Lemma \ref{lemma-minprimespolyhigher} instead of
Lemma \ref{lemma-minprimespoly} and
Lemma \ref{lemma-sysparhigher} instead of
Lemma \ref{lemma-syspar}.
We can use Lemma \ref{lemma-minprimespolyhigher} because $R$
is a catenary domain, so every height one prime ideal of $R$
has dimension $d - 1$, and hence the spectrum of $R/(f + h)$ is
equidimensional. For the convenience of the reader we write out
the details.

\medskip\noindent
Let $f \in J$ be a nonzero element.
We will modify $f$ slightly to obtain an element that generates a radical ideal.
The localization $R_\mathfrak p$ of $R$ at each height one prime
ideal $\mathfrak p$ is a discrete valuation ring, see
Algebra, Lemma \ref{algebra-lemma-characterize-dvr} or
Algebra, Lemma \ref{algebra-lemma-criterion-normal}.
We denote by $\text{ord}_\mathfrak p(f)$ the corresponding
valuation of $f$ in $R_{\mathfrak p}$. Let
$\mathfrak q_1, \ldots, \mathfrak q_s$
be the distinct height one prime ideals containing $f$.
Write $\text{ord}_{\mathfrak q_j}(f) = m_j \geq 1$
for each $j$. Then we define
$\text{div}(f) = \sum_{j = 1}^s m_j\mathfrak q_j$
as a formal linear combination of
height one primes with integer coefficients.
The ring $R/fR$ is reduced if and only if
$m_j = 1$ for $j = 1, \ldots, s$. Namely, if $m_j$ is $1$ then
$(R/fR)\mathfrak q_j$ is reduced and
$R/fR \subset \prod (R/fR)_{\mathfrak q_j}$ as
$\mathfrak q_1, \ldots, \mathfrak q_j$ are the associated primes
of $R/fR$, see Algebra, Lemmas
\ref{algebra-lemma-zero-at-ass-zero} and
\ref{algebra-lemma-normal-domain-intersection-localizations-height-1}.

\medskip\noindent
Choose and fix $g_2, \ldots, g_{d - 1}$ and $N$ as in
Lemma \ref{lemma-sysparhigher}.
For a nonzero $y \in R$ denote $t(y)$ the number of primes minimal over $y$.
Since $R$ is a normal domain, these primes
are height one and correspond $1$-to-$1$ to the minimal primes of
$R/yR$ (Algebra, Lemmas \ref{algebra-lemma-minimal-over-1} and
\ref{algebra-lemma-normal-domain-intersection-localizations-height-1}).
For example $t(f) = s$ is the number
of primes $\mathfrak q_j$ occurring in $\text{div}(f)$.
Let $h \in \mathfrak m^N$. Because $R$ is catenary, for each
height one prime $\mathfrak p$ of $R$ we have
$\dim(R/\mathfrak p) = d$. Hence by
Lemma \ref{lemma-minprimespolyhigher}
we have
\begin{align*}
t(f + h) & \leq \text{length}_{R/(f + h)}(R/(f + h, g_1, \ldots, g_{d - 1})) \\
& = \text{length}_R(R/(f + h, g_1, \ldots, g_{d - 1})) \\
& = \text{length}_R(R/(f, g_1, \ldots, g_{d - 1}))
\end{align*}
see Algebra, Lemma \ref{algebra-lemma-length-independent}
for the first equality.
Therefore we see that $t(f + h)$ is bounded independent of
$h \in \mathfrak m^N$.

\medskip\noindent
By the boundedness proved above we may pick $h \in \mathfrak m^N \cap J$
such that $t(f + h)$ is maximal among such $h$. Set $f' = f + h$.
Given $h' \in \mathfrak m^N \cap J$
we see that the number $t(f' + h') \leq t(f + h)$.
Thus after replacing $f$ by $f'$ we may assume that for every
$h \in \mathfrak m^N \cap J$ we have $t(f + h) \leq s$.

\medskip\noindent
Next, assume that we can find an element $h \in \mathfrak m^N \cap J$
such that for each $j$ we have $\text{ord}_{\mathfrak q_j}(h) \geq 1$ and
$\text{ord}_{\mathfrak q_j}(h) = 1 \Leftrightarrow m_j > 1$.
Then $\text{ord}_{\mathfrak q_j}(f + h) = 1$
for every $j$ by elementary properties of valuations.
Thus
$$
\text{div}(f + h) = \sum\nolimits_{j = 1}^s \mathfrak q_j +
\sum\nolimits_{k = 1}^v e_k \mathfrak r_k
$$
for some pairwise distinct height one prime ideals
$\mathfrak r_1, \ldots, \mathfrak r_v$ and $e_k \geq 1$.
However, since $s = t(f) \geq t(f + h)$ we see that $v = 0$
and we have found the desired element.

\medskip\noindent
Now we will pick $h$ that satisfies the above criteria.
By prime avoidance (Algebra, Lemma \ref{algebra-lemma-silly})
for each $1 \leq j \leq s$ we can find an element
$a_j \in \mathfrak q_j \cap J$
such that $a_j \not \in \mathfrak q_{j'}$ for $j' \not = j$.
Next, we can pick $b_j \in J \cap \mathfrak q_1 \cap \ldots \cap q_s$
with $b_j \not \in \mathfrak q_j^{(2)}$. Here
$\mathfrak q_j^{(2)} = \{x \in R \mid \text{ord}_{\mathfrak q_j}(x) \geq 2\}$
is the second symbolic power of $\mathfrak q_j$.
Prime avoidance applies because the ideal
$J' = J \cap \mathfrak q_1 \cap \ldots \cap q_s$
is radical, hence $R/J'$ is reduced, hence $(R/J')_{\mathfrak q_j}$
is reduced, hence $J'$ contains an element $x$ with
$\text{ord}_{\mathfrak q_j}(x) = 1$, hence
$J' \not \subset \mathfrak q_j^{(2)}$. Then the element
$$
c = \sum\nolimits_{j = 1, \ldots, s}
b_j \times \prod\nolimits_{j' \not = j} a_{j'}
$$
is an element of $J$
with $\text{ord}_{\mathfrak q_j}(c) = 1$ for all $j = 1, \ldots, s$
by elementary properties of valuations. Finally, we let
$$
h = c \times \prod\nolimits_{m_j = 1} a_j \times y
$$
where $y \in \mathfrak m^N$ is an element
which is not contained in $\mathfrak q_j$ for all $j$.
\end{proof}

\end{document}
